\documentclass[11pt,a4paper]{article}
\usepackage[utf8]{inputenc}
\usepackage[T1]{fontenc}
\usepackage[ngerman]{babel}
\usepackage{helvet}
\renewcommand{\familydefault}{\sfdefault}
\usepackage[left=2cm, right=2cm, top=2cm, bottom=2cm]{geometry}
\usepackage{amsmath,amssymb}
\usepackage{array}
\usepackage{booktabs}
\usepackage{tikz}
\usepackage{pgfplots}
\pgfplotsset{compat=1.18}
\usepackage{fancyhdr}
\usepackage{xcolor}

% Lösungsfarbe
\definecolor{correct}{RGB}{34,139,94}

\pagestyle{fancy}
\fancyhf{}
\fancyhead[L]{\small Physik Klasse 10E}
\fancyhead[R]{\small\bfseries ERWARTUNGSHORIZONT -- Variante B}
\fancyfoot[C]{\small Seite \thepage}
\renewcommand{\headrulewidth}{0.4pt}

\setlength{\parindent}{0pt}
\setlength{\fboxsep}{8pt}

% Aufgaben-Befehl im Layout B Stil
\newcounter{aufgabe}
\newcommand{\aufgabe}[2]{%
\stepcounter{aufgabe}%
\vspace{0.5cm}%
\noindent\fbox{\parbox{\dimexpr\textwidth-2\fboxsep-2\fboxrule}{%
\textbf{Aufgabe \theaufgabe: #1}\hfill\textbf{#2 BE}%
}}%
\vspace{0.3cm}%
}

\newcommand{\lsg}[1]{\textcolor{correct}{#1}}
\newcommand{\be}[1]{\hfill\textcolor{correct}{\small(#1 BE)}}

\begin{document}

% === KOPFBEREICH ===
\noindent\fbox{\parbox{\dimexpr\textwidth-2\fboxsep-2\fboxrule}{
\centering
{\Large\bfseries Probeklausur: Kinematik -- Variante B}\\[0.2cm]
{\large\bfseries Erwartungshorizont}\\[0.2cm]
{\small Klasse 10E \quad|\quad 90 Minuten \quad|\quad 60 BE}
}}

\vspace{0.4cm}

% === NOTENSPIEGEL ===
\begin{center}
\small
\begin{tabular}{|c|c|c|c|c|c|c|c|c|c|c|c|c|c|c|c|}
\hline
\textbf{NP} & 15 & 14 & 13 & 12 & 11 & 10 & 9 & 8 & 7 & 6 & 5 & 4 & 3 & 2 & 1 \\
\hline
\textbf{ab BE} & 60 & 59 & 58 & 55 & 51 & 48 & 44 & 41 & 36 & 32 & 29 & 24 & 20 & 17 & 12 \\
\hline
\end{tabular}

\vspace{0.2cm}
{\footnotesize NP = Notenpunkte (15-Punkte-Skala, Sek II MV)}
\end{center}

\vspace{0.3cm}

% === LÖSUNGEN ===

\aufgabe{Grundlagen}{6}

\begin{enumerate}
\item[a)] \lsg{90 km/h = 90 $\div$ 3,6 = 25 m/s} \be{1}

\item[b)] \lsg{Formel: $s = \frac{1}{2}gt^2$ $\Rightarrow$ $t = \sqrt{\frac{2s}{g}}$} \be{1}\\
\lsg{Einsetzen: $t = \sqrt{\frac{2 \cdot 20\,\mathrm{m}}{10\,\mathrm{m/s^2}}} = \sqrt{4\,\mathrm{s^2}} = 2\,\mathrm{s}$} \be{1}

\item[c)] Tabelle: \be{3}

\begin{center}
\begin{tabular}{|l|c|c|}
\hline
\textbf{Diagramm} & \textbf{Steigung bedeutet} & \textbf{Fläche bedeutet} \\
\hline
$s(t)$-Diagramm & \lsg{Geschwindigkeit $v$} & --- \\
\hline
$v(t)$-Diagramm & \lsg{Beschleunigung $a$} & \lsg{Weg $s$} \\
\hline
\end{tabular}
\end{center}
{\small\color{correct} Je 1 BE pro richtige Eintragung}
\end{enumerate}

\aufgabe{Notbremsung auf der Autobahn}{8}

\begin{enumerate}
\item[a)] \lsg{$v_0 = 90\,\mathrm{km/h} = 90 \div 3{,}6 = 25\,\mathrm{m/s}$} \be{1}

\item[b)] \lsg{Formel: $v = v_0 + a \cdot t$ $\Rightarrow$ $t = \dfrac{v_0}{|a|}$} \be{1}\\[0.2cm]
\lsg{Einsetzen: $t = \dfrac{25\,\mathrm{m/s}}{5\,\mathrm{m/s^2}} = 5\,\mathrm{s}$} \be{1}

\item[c)] \lsg{Formel: $s = \dfrac{v_0^2}{2|a|}$} \be{1}\\[0.2cm]
\lsg{Einsetzen: $s = \dfrac{(25\,\mathrm{m/s})^2}{2 \cdot 5\,\mathrm{m/s^2}} = \dfrac{625}{10}\,\mathrm{m} = 62{,}5\,\mathrm{m}$} \be{1}

\item[d)] \lsg{Nein, der LKW kann nicht rechtzeitig anhalten.} \be{1}\\
\lsg{Begründung: Der Bremsweg (62,5 m) ist länger als der Abstand zum Stauende (50 m).} \be{1}\\
\lsg{Der LKW fehlen 12,5 m: $62{,}5\,\mathrm{m} - 50\,\mathrm{m} = 12{,}5\,\mathrm{m}$} \be{1}
\end{enumerate}

\aufgabe{Fußball-Abstoß}{8}

\begin{enumerate}
\item[a)] \lsg{Formel: $v = v_0 - g \cdot t$ $\Rightarrow$ $t_{\text{steig}} = \dfrac{v_0}{g}$} \be{1}\\[0.2cm]
\lsg{Einsetzen: $t_{\text{steig}} = \dfrac{20\,\mathrm{m/s}}{10\,\mathrm{m/s^2}} = 2\,\mathrm{s}$} \be{1}

\item[b)] \lsg{Formel: $h_{\text{max}} = \dfrac{v_0^2}{2g}$} \be{1}\\[0.2cm]
\lsg{Einsetzen: $h_{\text{max}} = \dfrac{(20\,\mathrm{m/s})^2}{2 \cdot 10\,\mathrm{m/s^2}} = \dfrac{400}{20}\,\mathrm{m} = 20\,\mathrm{m}$} \be{2}

\item[c)] \lsg{Der Ball kommt mit $v = 20\,\mathrm{m/s}$ (nach unten) an.} \be{1}\\
\lsg{Begründung: Energieerhaltung / Symmetrie der Bewegung / zeitlose Formel liefert $|v| = v_0$} \be{2}
\end{enumerate}

\newpage

\aufgabe{Fahrradfahrt}{8}

\begin{enumerate}
\item[a)] \lsg{Phase A (0--6 s): Beschleunigte Bewegung, das Fahrrad beschleunigt von 0 auf 6 m/s.} \be{1}\\
\lsg{Phase B (6--12 s): Gleichförmige Bewegung mit konstant 6 m/s.} \be{1}\\
\lsg{Phase C (12--15 s): Bremsvorgang, das Fahrrad bremst von 6 m/s auf 0 ab.} \be{1}

\item[b)] \lsg{Formel: $a = \dfrac{\Delta v}{\Delta t}$} \be{1}\\[0.2cm]
\lsg{Einsetzen: $a = \dfrac{6\,\mathrm{m/s} - 0}{6\,\mathrm{s}} = 1\,\mathrm{m/s^2}$} \be{1}

\item[c)] \lsg{Fläche unter dem Graphen = Weg. Zerlegung in geometrische Figuren:} \be{1}\\
\lsg{Phase A (Dreieck): $s_A = \dfrac{1}{2} \cdot 6\,\mathrm{s} \cdot 6\,\mathrm{m/s} = 18\,\mathrm{m}$}\\
\lsg{Phase B (Rechteck): $s_B = 6\,\mathrm{s} \cdot 6\,\mathrm{m/s} = 36\,\mathrm{m}$}\\
\lsg{Phase C (Dreieck): $s_C = \dfrac{1}{2} \cdot 3\,\mathrm{s} \cdot 6\,\mathrm{m/s} = 9\,\mathrm{m}$}\\
\lsg{Gesamtweg: $s = 18 + 36 + 9 = 63\,\mathrm{m}$} \be{2}
\end{enumerate}

\aufgabe{Fallschirm und Luftwiderstand}{8}

\begin{enumerate}
\item[a)] \lsg{Bei zunehmender Geschwindigkeit steigt der Luftwiderstand ($F_W \sim v^2$).} \be{1}\\
\lsg{Wenn $F_W = F_G$ (Gewichtskraft), heben sich die Kräfte auf.} \be{1}\\
\lsg{Der Springer fällt dann mit konstanter Grenzgeschwindigkeit.} \be{1}

\item[b)] \lsg{Beide Kugeln haben die gleiche Querschnittsfläche $A$ (gleiche Größe).} \be{1}\\
\lsg{Die Stahlkugel hat mehr Masse, also größere Gewichtskraft $F_G$.} \be{1}\\
\lsg{Daher erreicht sie eine höhere Grenzgeschwindigkeit und kommt schneller an.} \be{1}

\item[c)] \lsg{Die Aussage ist \textbf{richtig}. Das Experiment auf dem Mond beweist es.} \be{1}\\
\lsg{Im Vakuum gibt es keinen Luftwiderstand. Die Formel $s = \frac{1}{2}gt^2$ enthält keine Masse.} \be{1}
\end{enumerate}

\aufgabe{Überholvorgang auf der Landstraße}{10}

\begin{enumerate}
\item[a)] \lsg{$v_{\text{PKW}} = 72\,\mathrm{km/h} = 20\,\mathrm{m/s}$}\\
\lsg{$v_{\text{LKW}} = 54\,\mathrm{km/h} = 15\,\mathrm{m/s}$} \be{2}

\item[b)] \lsg{Relativgeschwindigkeit: $v_{\text{rel}} = v_{\text{PKW}} - v_{\text{LKW}} = 20 - 15 = 5\,\mathrm{m/s}$} \be{2}

\item[c)] \lsg{Zeit: $t = \dfrac{s}{v_{\text{rel}}} = \dfrac{20\,\mathrm{m}}{5\,\mathrm{m/s}} = 4\,\mathrm{s}$} \be{2}

\item[d)] \lsg{Strecke: $s = v \cdot t = 20\,\mathrm{m/s} \cdot 4\,\mathrm{s} = 80\,\mathrm{m}$} \be{2}

\item[e)] \lsg{Bei geringem Geschwindigkeitsunterschied dauert der Überholvorgang lange.}\\
\lsg{Man ist lange auf der Gegenfahrbahn, was bei Gegenverkehr gefährlich ist.} \be{2}
\end{enumerate}

\newpage

\aufgabe{Sprung vom 10-Meter-Turm}{10}

\begin{enumerate}
\item[a)] \lsg{Formel: $s = \frac{1}{2}gt^2$ $\Rightarrow$ $t = \sqrt{\frac{2s}{g}} = \sqrt{\frac{20}{10}} = \sqrt{2} \approx 1{,}41\,\mathrm{s}$} \be{2}

\item[b)] \lsg{$v = g \cdot t = 10 \cdot \sqrt{2} \approx 14{,}1\,\mathrm{m/s}$}\\
\lsg{Alternativ: $v = \sqrt{2gs} = \sqrt{200} \approx 14{,}1\,\mathrm{m/s}$} \be{2}

\item[c)] \lsg{Zeitlose Formel: $v^2 = v_0^2 + 2gs$}\\
\lsg{$v^2 = (2\,\mathrm{m/s})^2 + 2 \cdot 10\,\mathrm{m/s^2} \cdot 10\,\mathrm{m} = 4 + 200 = 204\,\mathrm{m^2/s^2}$}\\
\lsg{$v = \sqrt{204} \approx 14{,}3\,\mathrm{m/s}$} \be{3}

\item[d)] \lsg{Der Unterschied beträgt nur etwa 0,2 m/s (14,3 vs. 14,1 m/s).}\\
\lsg{Die zusätzliche kinetische Energie durch $v_0 = 2\,\mathrm{m/s}$ ist gering im Vergleich zur potentiellen Energie aus 10 m Höhe.} \be{3}
\end{enumerate}

\vspace{0.3cm}
\hrule
\vspace{0.3cm}

\noindent\fbox{\parbox{\dimexpr\textwidth-2\fboxsep-2\fboxrule}{%
\textbf{Zusatzaufgabe} (max. 2 Zusatz-BE)
}}

\vspace{0.3cm}

\lsg{Ansatz: Positionen beider Fahrzeuge als Funktion der Zeit aufstellen.}\\[0.2cm]
\lsg{Auto: $s_{\text{Auto}}(t) = \frac{1}{2} \cdot 2\,\mathrm{m/s^2} \cdot t^2 = t^2$}\\
\lsg{Radfahrer (startet 25 m dahinter): $s_{\text{Rad}}(t) = 10\,\mathrm{m/s} \cdot t - 25\,\mathrm{m}$}\\[0.2cm]
\lsg{Gleichsetzen: $t^2 = 10t - 25$}\\
\lsg{$t^2 - 10t + 25 = 0$}\\
\lsg{$(t - 5)^2 = 0$}\\
\lsg{$t = 5\,\mathrm{s}$} \be{2}

{\small\color{correct} Probe: $s_{\text{Auto}}(5) = 25\,\mathrm{m}$, $s_{\text{Rad}}(5) = 25\,\mathrm{m}$ \checkmark}

\vspace{0.3cm}
\hrule
\vspace{0.3cm}

\begin{center}
\fbox{\parbox{0.8\textwidth}{
\centering
\textbf{Zusammenfassung der Anforderungsbereiche}\\[0.3cm]
\begin{tabular}{|l|c|c|c|}
\hline
\textbf{AFB} & \textbf{Aufgaben} & \textbf{BE} & \textbf{Anteil} \\
\hline
I (Reproduktion) & 1 & 6 & 10\% \\
\hline
II (Anwendung) & 2, 3, 4, 6 & 34 & 59\% \\
\hline
III (Transfer) & 5, 7, Zusatz & 18 (+2) & 31\% \\
\hline
\textbf{Gesamt} & & \textbf{58 (+2)} & 100\% \\
\hline
\end{tabular}
}}
\end{center}

\end{document}
